
\documentclass{article}
\usepackage{colortbl}
\usepackage{makecell}
\usepackage{multirow}
\usepackage{supertabular}

\begin{document}

\newcounter{utterance}

\centering \large Interaction Transcript for game `codenames', experiment `risk\_low', episode 1 with gpt{-}5.4{-}azure.
\vspace{24pt}

{ \footnotesize  \setcounter{utterance}{1}
\setlength{\tabcolsep}{0pt}
\begin{supertabular}{c@{$\;$}|p{.15\linewidth}@{}p{.15\linewidth}p{.15\linewidth}p{.15\linewidth}p{.15\linewidth}p{.15\linewidth}}
    \# & $\;$A & \multicolumn{4}{c}{Game Master} & $\;\:$B\\
    \hline

    \theutterance \stepcounter{utterance}  
    & & & \multicolumn{2}{p{0.3\linewidth}}{
        \cellcolor[rgb]{0.9,0.9,0.9}{
            \makecell[{{p{\linewidth}}}]{
                \texttt{\tiny{[GM$|$GM]}}
                {'hidden': {'team': ['pantalon', 'cour', 'eau', 'stock', 'araignée', 'laser', 'papier', 'flûte', 'serrure'], 'innocent': ['arc', 'traverser', 'carré', 'taupe', 'scorpion', 'police', 'triangle', 'avion'], 'opponent': ['bâton', 'mercure', 'pirate', 'tuyau', 'âme', 'sorcière', 'balle', 'balançoire'], 'assassin': []}, 'revealed': {'team': {'team': [], 'innocent': [], 'opponent': [], 'assassin': []}, 'opponent': {'team': [], 'innocent': [], 'opponent': [], 'assassin': []}}}
            }
        }
    }
    & & \\ \\

    \theutterance \stepcounter{utterance}  
    & & \multicolumn{4}{p{0.6\linewidth}}{
        \cellcolor[rgb]{0.9,0.9,0.9}{
            \makecell[{{p{\linewidth}}}]{
                \texttt{\tiny{[P1$\langle$GM]}}
                \texttt{Fournissez un indice d'un seul mot qui se rapporte à certains des mots de votre équipe figurant dans la liste ci{-}dessous. Vous pouvez choisir à combien et à quels mots votre indice se rapporte.} \\
\texttt{Le but est que votre partenaire devine tous les mots de votre équipe à partir des indices que vous lui donnez, avant que l'autre équipe n'ait découvert tous les siens. Si votre partenaire devine un mot « adversaire » ou « distracteur », son tour de devinette s'achève et l'équipe adverse peut jouer à son tour. Si votre partenaire devine un mot « mortel », le jeu s'arrête et votre équipe perd immédiatement.} \\
\texttt{Plus votre indice se rapporte à de mots, plus vous pouvez gagner rapidement le jeu, mais plus les associations peuvent être difficiles à deviner pour votre partenaire, alors choisissez bien votre mot indice. Le mot indice doit être sémantiquement lié aux mots cibles, il ne peut pas être l'un des mots des listes ou en contenir des parties.} \\
\\ 
\texttt{Donnez toujours votre indice d'un seul mot et votre liste de mots cibles séparés par des virgules dans le format suivant et faites en sorte que vos réponses soient aussi courtes que possible, n'incluez jamais d'autre texte que celui requis dans ce format :} \\
\\ 
\texttt{INDICE: <MOT>} \\
\texttt{CIBLES: <CIBLES>} \\
\\ 
\texttt{Les mots de votre équipe sont : pantalon, cour, eau, stock, araignée, laser, papier, flûte, serrure.} \\
\texttt{Les mots de l'équipe adverse sont : bâton, mercure, pirate, tuyau, âme, sorcière, balle, balançoire.} \\
\texttt{Les mots distracteurs sont : arc, traverser, carré, taupe, scorpion, police, triangle, avion.} \\
\texttt{Les mots mortels sont : .} \\
            }
        }
    }
    & \\ \\

    \theutterance \stepcounter{utterance}  
    & \multicolumn{4}{p{0.6\linewidth}}{
        \cellcolor[rgb]{0.9,0.9,0.9}{
            \makecell[{{p{\linewidth}}}]{
                \texttt{\tiny{[P1$\rangle$GM]}}
                \texttt{INDICE: imprimante} \\
\texttt{CIBLES: laser, papier} \\
            }
        }
    }
    & & \\ \\

    \theutterance \stepcounter{utterance}  
    & & & \multicolumn{2}{p{0.3\linewidth}}{
        \cellcolor[rgb]{0.9,0.9,0.9}{
            \makecell[{{p{\linewidth}}}]{
                \texttt{\tiny{[GM$|$GM]}}
                \texttt{imprimante} \\
            }
        }
    }
    & & \\ \\

    \theutterance \stepcounter{utterance}  
    & & & \multicolumn{2}{p{0.3\linewidth}}{
        \cellcolor[rgb]{0.9,0.9,0.9}{
            \makecell[{{p{\linewidth}}}]{
                \texttt{\tiny{[GM$|$GM]}}
                ['laser', 'papier']
            }
        }
    }
    & & \\ \\

    \theutterance \stepcounter{utterance}  
    & & & \multicolumn{2}{p{0.3\linewidth}}{
        \cellcolor[rgb]{0.9,0.9,0.9}{
            \makecell[{{p{\linewidth}}}]{
                \texttt{\tiny{[GM$|$GM]}}
                {'word': 'laser', 'assignment': 'team'}
            }
        }
    }
    & & \\ \\

    \theutterance \stepcounter{utterance}  
    & & & \multicolumn{2}{p{0.3\linewidth}}{
        \cellcolor[rgb]{0.9,0.9,0.9}{
            \makecell[{{p{\linewidth}}}]{
                \texttt{\tiny{[GM$|$GM]}}
                {'word': 'papier', 'assignment': 'team'}
            }
        }
    }
    & & \\ \\

    \theutterance \stepcounter{utterance}  
    & & \multicolumn{4}{p{0.6\linewidth}}{
        \cellcolor[rgb]{0.9,0.9,0.9}{
            \makecell[{{p{\linewidth}}}]{
                \texttt{\tiny{[GM$\rangle$P2]}}
                \texttt{Fournissez une liste séparée par des virgules d'un maximum de 2 mots de la liste suivante qui se rapportent le mieux ou sont le plus étroitement associés au mot "imprimante". Commencez toujours votre liste de supposition(s) par DEVINETTE:  et n'incluez aucun autre texte dans votre réponse.} \\
\\ 
\texttt{serrure, sorcière, araignée, pirate, taupe, balançoire, police, âme, bâton, scorpion, pantalon, avion, stock, mercure, laser, tuyau, balle, eau, traverser, triangle, arc, papier, carré, flûte, cour} \\
            }
        }
    }
    & \\ \\

    \theutterance \stepcounter{utterance}  
    & & & \multicolumn{4}{p{0.6\linewidth}}{
        \cellcolor[rgb]{0.9,0.9,0.9}{
            \makecell[{{p{\linewidth}}}]{
                \texttt{\tiny{[GM$\langle$P2]}}
                \texttt{DEVINETTE: papier, laser} \\
            }
        }
    }
     \\ \\

    \theutterance \stepcounter{utterance}  
    & & & \multicolumn{2}{p{0.3\linewidth}}{
        \cellcolor[rgb]{0.9,0.9,0.9}{
            \makecell[{{p{\linewidth}}}]{
                \texttt{\tiny{[GM$|$GM]}}
                ['papier', 'laser']
            }
        }
    }
    & & \\ \\

    \theutterance \stepcounter{utterance}  
    & & & \multicolumn{2}{p{0.3\linewidth}}{
        \cellcolor[rgb]{0.9,0.9,0.9}{
            \makecell[{{p{\linewidth}}}]{
                \texttt{\tiny{[GM$|$GM]}}
                {'word': 'papier', 'assignment': 'team'}
            }
        }
    }
    & & \\ \\

    \theutterance \stepcounter{utterance}  
    & & & \multicolumn{2}{p{0.3\linewidth}}{
        \cellcolor[rgb]{0.9,0.9,0.9}{
            \makecell[{{p{\linewidth}}}]{
                \texttt{\tiny{[GM$|$GM]}}
                {'word': 'papier', 'assignment': 'team'}
            }
        }
    }
    & & \\ \\

    \theutterance \stepcounter{utterance}  
    & & & \multicolumn{2}{p{0.3\linewidth}}{
        \cellcolor[rgb]{0.9,0.9,0.9}{
            \makecell[{{p{\linewidth}}}]{
                \texttt{\tiny{[GM$|$GM]}}
                {'word': 'laser', 'assignment': 'team'}
            }
        }
    }
    & & \\ \\

    \theutterance \stepcounter{utterance}  
    & & & \multicolumn{2}{p{0.3\linewidth}}{
        \cellcolor[rgb]{0.9,0.9,0.9}{
            \makecell[{{p{\linewidth}}}]{
                \texttt{\tiny{[GM$|$GM]}}
                {'word': 'laser', 'assignment': 'team'}
            }
        }
    }
    & & \\ \\

    \theutterance \stepcounter{utterance}  
    & & & \multicolumn{2}{p{0.3\linewidth}}{
        \cellcolor[rgb]{0.9,0.9,0.9}{
            \makecell[{{p{\linewidth}}}]{
                \texttt{\tiny{[GM$|$GM]}}
                {'word': 'mercure', 'assignment': 'opponent'}
            }
        }
    }
    & & \\ \\

    \theutterance \stepcounter{utterance}  
    & & & \multicolumn{2}{p{0.3\linewidth}}{
        \cellcolor[rgb]{0.9,0.9,0.9}{
            \makecell[{{p{\linewidth}}}]{
                \texttt{\tiny{[GM$|$GM]}}
                {'hidden': {'team': ['pantalon', 'cour', 'eau', 'stock', 'araignée', 'flûte', 'serrure'], 'innocent': ['arc', 'traverser', 'carré', 'taupe', 'scorpion', 'police', 'triangle', 'avion'], 'opponent': ['bâton', 'pirate', 'tuyau', 'âme', 'sorcière', 'balle', 'balançoire'], 'assassin': []}, 'revealed': {'team': {'team': ['papier', 'laser'], 'innocent': [], 'opponent': [], 'assassin': []}, 'opponent': {'team': [], 'innocent': [], 'opponent': ['mercure'], 'assassin': []}}}
            }
        }
    }
    & & \\ \\

    \theutterance \stepcounter{utterance}  
    & & \multicolumn{4}{p{0.6\linewidth}}{
        \cellcolor[rgb]{0.9,0.9,0.9}{
            \makecell[{{p{\linewidth}}}]{
                \texttt{\tiny{[P1$\langle$GM]}}
                \texttt{Les mots papier, laser ont été devinés correctement. Le tour de votre coéquipier s'est terminé là.} \\
\texttt{Fournissez maintenant un autre indice relatif à certains des mots restants de votre équipe et une liste des mots cibles correspondants. N'oubliez pas de commencer votre indice par 'INDICE: ', de mettre une nouvelle ligne, et de commencer votre liste de mots cibles séparés par des virgules par 'CIBLES: '. Remarque : certains mots ont été supprimés des listes par rapport aux demandes précédentes.} \\
\\ 
\texttt{Les mots restants de votre équipe sont : pantalon, cour, eau, stock, araignée, flûte, serrure.} \\
\texttt{Les mots restants pour votre adversaire sont : bâton, mercure, pirate, tuyau, âme, sorcière, balle, balançoire.} \\
\texttt{Les mots distracteurs restants sont : arc, traverser, carré, taupe, scorpion, police, triangle, avion.} \\
\texttt{Les mots tueurs restants sont : .} \\
            }
        }
    }
    & \\ \\

    \theutterance \stepcounter{utterance}  
    & \multicolumn{4}{p{0.6\linewidth}}{
        \cellcolor[rgb]{0.9,0.9,0.9}{
            \makecell[{{p{\linewidth}}}]{
                \texttt{\tiny{[P1$\rangle$GM]}}
                \texttt{INDICE: cadenas} \\
\texttt{CIBLES: serrure} \\
            }
        }
    }
    & & \\ \\

    \theutterance \stepcounter{utterance}  
    & & & \multicolumn{2}{p{0.3\linewidth}}{
        \cellcolor[rgb]{0.9,0.9,0.9}{
            \makecell[{{p{\linewidth}}}]{
                \texttt{\tiny{[GM$|$GM]}}
                \texttt{cadenas} \\
            }
        }
    }
    & & \\ \\

    \theutterance \stepcounter{utterance}  
    & & & \multicolumn{2}{p{0.3\linewidth}}{
        \cellcolor[rgb]{0.9,0.9,0.9}{
            \makecell[{{p{\linewidth}}}]{
                \texttt{\tiny{[GM$|$GM]}}
                ['serrure']
            }
        }
    }
    & & \\ \\

    \theutterance \stepcounter{utterance}  
    & & & \multicolumn{2}{p{0.3\linewidth}}{
        \cellcolor[rgb]{0.9,0.9,0.9}{
            \makecell[{{p{\linewidth}}}]{
                \texttt{\tiny{[GM$|$GM]}}
                {'word': 'serrure', 'assignment': 'team'}
            }
        }
    }
    & & \\ \\

    \theutterance \stepcounter{utterance}  
    & & \multicolumn{4}{p{0.6\linewidth}}{
        \cellcolor[rgb]{0.9,0.9,0.9}{
            \makecell[{{p{\linewidth}}}]{
                \texttt{\tiny{[GM$\rangle$P2]}}
                \texttt{Les mots papier, laser ont été devinés correctement. Votre tour s'est terminé là.} \\
\texttt{Fournissez maintenant une autre liste, séparée par des virgules, d'au moins 1 et jusqu'à 1 mots de la liste suivante de mots qui se rapportent le mieux ou sont le plus étroitement associés au mot 'cadenas'. N'oubliez pas de commencer votre réponse par DEVINETTE: . Remarque : certains mots ont été supprimés de la liste par rapport aux demandes précédentes.} \\
\\ 
\texttt{serrure, sorcière, araignée, pirate, taupe, balançoire, police, âme, bâton, scorpion, pantalon, avion, stock, tuyau, balle, eau, traverser, triangle, arc, carré, flûte, cour} \\
            }
        }
    }
    & \\ \\

    \theutterance \stepcounter{utterance}  
    & & & \multicolumn{4}{p{0.6\linewidth}}{
        \cellcolor[rgb]{0.9,0.9,0.9}{
            \makecell[{{p{\linewidth}}}]{
                \texttt{\tiny{[GM$\langle$P2]}}
                \texttt{DEVINETTE: serrure} \\
            }
        }
    }
     \\ \\

    \theutterance \stepcounter{utterance}  
    & & & \multicolumn{2}{p{0.3\linewidth}}{
        \cellcolor[rgb]{0.9,0.9,0.9}{
            \makecell[{{p{\linewidth}}}]{
                \texttt{\tiny{[GM$|$GM]}}
                ['serrure']
            }
        }
    }
    & & \\ \\

    \theutterance \stepcounter{utterance}  
    & & & \multicolumn{2}{p{0.3\linewidth}}{
        \cellcolor[rgb]{0.9,0.9,0.9}{
            \makecell[{{p{\linewidth}}}]{
                \texttt{\tiny{[GM$|$GM]}}
                {'word': 'serrure', 'assignment': 'team'}
            }
        }
    }
    & & \\ \\

    \theutterance \stepcounter{utterance}  
    & & & \multicolumn{2}{p{0.3\linewidth}}{
        \cellcolor[rgb]{0.9,0.9,0.9}{
            \makecell[{{p{\linewidth}}}]{
                \texttt{\tiny{[GM$|$GM]}}
                {'word': 'serrure', 'assignment': 'team'}
            }
        }
    }
    & & \\ \\

    \theutterance \stepcounter{utterance}  
    & & & \multicolumn{2}{p{0.3\linewidth}}{
        \cellcolor[rgb]{0.9,0.9,0.9}{
            \makecell[{{p{\linewidth}}}]{
                \texttt{\tiny{[GM$|$GM]}}
                {'word': 'balle', 'assignment': 'opponent'}
            }
        }
    }
    & & \\ \\

    \theutterance \stepcounter{utterance}  
    & & & \multicolumn{2}{p{0.3\linewidth}}{
        \cellcolor[rgb]{0.9,0.9,0.9}{
            \makecell[{{p{\linewidth}}}]{
                \texttt{\tiny{[GM$|$GM]}}
                {'hidden': {'team': ['pantalon', 'cour', 'eau', 'stock', 'araignée', 'flûte'], 'innocent': ['arc', 'traverser', 'carré', 'taupe', 'scorpion', 'police', 'triangle', 'avion'], 'opponent': ['bâton', 'pirate', 'tuyau', 'âme', 'sorcière', 'balançoire'], 'assassin': []}, 'revealed': {'team': {'team': ['papier', 'laser', 'serrure'], 'innocent': [], 'opponent': [], 'assassin': []}, 'opponent': {'team': [], 'innocent': [], 'opponent': ['mercure', 'balle'], 'assassin': []}}}
            }
        }
    }
    & & \\ \\

    \theutterance \stepcounter{utterance}  
    & & \multicolumn{4}{p{0.6\linewidth}}{
        \cellcolor[rgb]{0.9,0.9,0.9}{
            \makecell[{{p{\linewidth}}}]{
                \texttt{\tiny{[P1$\langle$GM]}}
                \texttt{Le mot serrure a été deviné correctement. Le tour de votre coéquipier s'est terminé là.} \\
\texttt{Fournissez maintenant un autre indice relatif à certains des mots restants de votre équipe et une liste des mots cibles correspondants. N'oubliez pas de commencer votre indice par 'INDICE: ', de mettre une nouvelle ligne, et de commencer votre liste de mots cibles séparés par des virgules par 'CIBLES: '. Remarque : certains mots ont été supprimés des listes par rapport aux demandes précédentes.} \\
\\ 
\texttt{Les mots restants de votre équipe sont : pantalon, cour, eau, stock, araignée, flûte.} \\
\texttt{Les mots restants pour votre adversaire sont : bâton, pirate, tuyau, âme, sorcière, balle, balançoire.} \\
\texttt{Les mots distracteurs restants sont : arc, traverser, carré, taupe, scorpion, police, triangle, avion.} \\
\texttt{Les mots tueurs restants sont : .} \\
            }
        }
    }
    & \\ \\

    \theutterance \stepcounter{utterance}  
    & \multicolumn{4}{p{0.6\linewidth}}{
        \cellcolor[rgb]{0.9,0.9,0.9}{
            \makecell[{{p{\linewidth}}}]{
                \texttt{\tiny{[P1$\rangle$GM]}}
                \texttt{INDICE: toile} \\
\texttt{CIBLES: araignée} \\
            }
        }
    }
    & & \\ \\

    \theutterance \stepcounter{utterance}  
    & & & \multicolumn{2}{p{0.3\linewidth}}{
        \cellcolor[rgb]{0.9,0.9,0.9}{
            \makecell[{{p{\linewidth}}}]{
                \texttt{\tiny{[GM$|$GM]}}
                \texttt{toile} \\
            }
        }
    }
    & & \\ \\

    \theutterance \stepcounter{utterance}  
    & & & \multicolumn{2}{p{0.3\linewidth}}{
        \cellcolor[rgb]{0.9,0.9,0.9}{
            \makecell[{{p{\linewidth}}}]{
                \texttt{\tiny{[GM$|$GM]}}
                ['araignée']
            }
        }
    }
    & & \\ \\

    \theutterance \stepcounter{utterance}  
    & & & \multicolumn{2}{p{0.3\linewidth}}{
        \cellcolor[rgb]{0.9,0.9,0.9}{
            \makecell[{{p{\linewidth}}}]{
                \texttt{\tiny{[GM$|$GM]}}
                {'word': 'araignée', 'assignment': 'team'}
            }
        }
    }
    & & \\ \\

    \theutterance \stepcounter{utterance}  
    & & \multicolumn{4}{p{0.6\linewidth}}{
        \cellcolor[rgb]{0.9,0.9,0.9}{
            \makecell[{{p{\linewidth}}}]{
                \texttt{\tiny{[GM$\rangle$P2]}}
                \texttt{Le mot serrure a été deviné correctement. Votre tour s'est terminé là.} \\
\texttt{Fournissez maintenant une autre liste, séparée par des virgules, d'au moins 1 et jusqu'à 1 mots de la liste suivante de mots qui se rapportent le mieux ou sont le plus étroitement associés au mot 'toile'. N'oubliez pas de commencer votre réponse par DEVINETTE: . Remarque : certains mots ont été supprimés de la liste par rapport aux demandes précédentes.} \\
\\ 
\texttt{sorcière, araignée, pirate, taupe, balançoire, police, âme, bâton, scorpion, pantalon, avion, stock, tuyau, eau, traverser, triangle, arc, carré, flûte, cour} \\
            }
        }
    }
    & \\ \\

    \theutterance \stepcounter{utterance}  
    & & & \multicolumn{4}{p{0.6\linewidth}}{
        \cellcolor[rgb]{0.9,0.9,0.9}{
            \makecell[{{p{\linewidth}}}]{
                \texttt{\tiny{[GM$\langle$P2]}}
                \texttt{DEVINETTE: araignée} \\
            }
        }
    }
     \\ \\

    \theutterance \stepcounter{utterance}  
    & & & \multicolumn{2}{p{0.3\linewidth}}{
        \cellcolor[rgb]{0.9,0.9,0.9}{
            \makecell[{{p{\linewidth}}}]{
                \texttt{\tiny{[GM$|$GM]}}
                ['araignée']
            }
        }
    }
    & & \\ \\

    \theutterance \stepcounter{utterance}  
    & & & \multicolumn{2}{p{0.3\linewidth}}{
        \cellcolor[rgb]{0.9,0.9,0.9}{
            \makecell[{{p{\linewidth}}}]{
                \texttt{\tiny{[GM$|$GM]}}
                {'word': 'araignée', 'assignment': 'team'}
            }
        }
    }
    & & \\ \\

    \theutterance \stepcounter{utterance}  
    & & & \multicolumn{2}{p{0.3\linewidth}}{
        \cellcolor[rgb]{0.9,0.9,0.9}{
            \makecell[{{p{\linewidth}}}]{
                \texttt{\tiny{[GM$|$GM]}}
                {'word': 'araignée', 'assignment': 'team'}
            }
        }
    }
    & & \\ \\

    \theutterance \stepcounter{utterance}  
    & & & \multicolumn{2}{p{0.3\linewidth}}{
        \cellcolor[rgb]{0.9,0.9,0.9}{
            \makecell[{{p{\linewidth}}}]{
                \texttt{\tiny{[GM$|$GM]}}
                {'word': 'balançoire', 'assignment': 'opponent'}
            }
        }
    }
    & & \\ \\

    \theutterance \stepcounter{utterance}  
    & & & \multicolumn{2}{p{0.3\linewidth}}{
        \cellcolor[rgb]{0.9,0.9,0.9}{
            \makecell[{{p{\linewidth}}}]{
                \texttt{\tiny{[GM$|$GM]}}
                {'hidden': {'team': ['pantalon', 'cour', 'eau', 'stock', 'flûte'], 'innocent': ['arc', 'traverser', 'carré', 'taupe', 'scorpion', 'police', 'triangle', 'avion'], 'opponent': ['bâton', 'pirate', 'tuyau', 'âme', 'sorcière'], 'assassin': []}, 'revealed': {'team': {'team': ['papier', 'laser', 'serrure', 'araignée'], 'innocent': [], 'opponent': [], 'assassin': []}, 'opponent': {'team': [], 'innocent': [], 'opponent': ['mercure', 'balle', 'balançoire'], 'assassin': []}}}
            }
        }
    }
    & & \\ \\

    \theutterance \stepcounter{utterance}  
    & & \multicolumn{4}{p{0.6\linewidth}}{
        \cellcolor[rgb]{0.9,0.9,0.9}{
            \makecell[{{p{\linewidth}}}]{
                \texttt{\tiny{[P1$\langle$GM]}}
                \texttt{Le mot araignée a été deviné correctement. Le tour de votre coéquipier s'est terminé là.} \\
\texttt{Fournissez maintenant un autre indice relatif à certains des mots restants de votre équipe et une liste des mots cibles correspondants. N'oubliez pas de commencer votre indice par 'INDICE: ', de mettre une nouvelle ligne, et de commencer votre liste de mots cibles séparés par des virgules par 'CIBLES: '. Remarque : certains mots ont été supprimés des listes par rapport aux demandes précédentes.} \\
\\ 
\texttt{Les mots restants de votre équipe sont : pantalon, cour, eau, stock, flûte.} \\
\texttt{Les mots restants pour votre adversaire sont : bâton, pirate, tuyau, âme, sorcière, balançoire.} \\
\texttt{Les mots distracteurs restants sont : arc, traverser, carré, taupe, scorpion, police, triangle, avion.} \\
\texttt{Les mots tueurs restants sont : .} \\
            }
        }
    }
    & \\ \\

    \theutterance \stepcounter{utterance}  
    & \multicolumn{4}{p{0.6\linewidth}}{
        \cellcolor[rgb]{0.9,0.9,0.9}{
            \makecell[{{p{\linewidth}}}]{
                \texttt{\tiny{[P1$\rangle$GM]}}
                \texttt{INDICE: piccolo} \\
\texttt{CIBLES: flûte} \\
            }
        }
    }
    & & \\ \\

    \theutterance \stepcounter{utterance}  
    & & & \multicolumn{2}{p{0.3\linewidth}}{
        \cellcolor[rgb]{0.9,0.9,0.9}{
            \makecell[{{p{\linewidth}}}]{
                \texttt{\tiny{[GM$|$GM]}}
                \texttt{piccolo} \\
            }
        }
    }
    & & \\ \\

    \theutterance \stepcounter{utterance}  
    & & & \multicolumn{2}{p{0.3\linewidth}}{
        \cellcolor[rgb]{0.9,0.9,0.9}{
            \makecell[{{p{\linewidth}}}]{
                \texttt{\tiny{[GM$|$GM]}}
                ['flûte']
            }
        }
    }
    & & \\ \\

    \theutterance \stepcounter{utterance}  
    & & & \multicolumn{2}{p{0.3\linewidth}}{
        \cellcolor[rgb]{0.9,0.9,0.9}{
            \makecell[{{p{\linewidth}}}]{
                \texttt{\tiny{[GM$|$GM]}}
                {'word': 'flûte', 'assignment': 'team'}
            }
        }
    }
    & & \\ \\

    \theutterance \stepcounter{utterance}  
    & & \multicolumn{4}{p{0.6\linewidth}}{
        \cellcolor[rgb]{0.9,0.9,0.9}{
            \makecell[{{p{\linewidth}}}]{
                \texttt{\tiny{[GM$\rangle$P2]}}
                \texttt{Le mot araignée a été deviné correctement. Votre tour s'est terminé là.} \\
\texttt{Fournissez maintenant une autre liste, séparée par des virgules, d'au moins 1 et jusqu'à 1 mots de la liste suivante de mots qui se rapportent le mieux ou sont le plus étroitement associés au mot 'piccolo'. N'oubliez pas de commencer votre réponse par DEVINETTE: . Remarque : certains mots ont été supprimés de la liste par rapport aux demandes précédentes.} \\
\\ 
\texttt{sorcière, pirate, taupe, police, âme, bâton, scorpion, pantalon, avion, stock, tuyau, eau, traverser, triangle, arc, carré, flûte, cour} \\
            }
        }
    }
    & \\ \\

    \theutterance \stepcounter{utterance}  
    & & & \multicolumn{4}{p{0.6\linewidth}}{
        \cellcolor[rgb]{0.9,0.9,0.9}{
            \makecell[{{p{\linewidth}}}]{
                \texttt{\tiny{[GM$\langle$P2]}}
                \texttt{DEVINETTE: flûte} \\
            }
        }
    }
     \\ \\

    \theutterance \stepcounter{utterance}  
    & & & \multicolumn{2}{p{0.3\linewidth}}{
        \cellcolor[rgb]{0.9,0.9,0.9}{
            \makecell[{{p{\linewidth}}}]{
                \texttt{\tiny{[GM$|$GM]}}
                ['flûte']
            }
        }
    }
    & & \\ \\

    \theutterance \stepcounter{utterance}  
    & & & \multicolumn{2}{p{0.3\linewidth}}{
        \cellcolor[rgb]{0.9,0.9,0.9}{
            \makecell[{{p{\linewidth}}}]{
                \texttt{\tiny{[GM$|$GM]}}
                {'word': 'flûte', 'assignment': 'team'}
            }
        }
    }
    & & \\ \\

    \theutterance \stepcounter{utterance}  
    & & & \multicolumn{2}{p{0.3\linewidth}}{
        \cellcolor[rgb]{0.9,0.9,0.9}{
            \makecell[{{p{\linewidth}}}]{
                \texttt{\tiny{[GM$|$GM]}}
                {'word': 'flûte', 'assignment': 'team'}
            }
        }
    }
    & & \\ \\

    \theutterance \stepcounter{utterance}  
    & & & \multicolumn{2}{p{0.3\linewidth}}{
        \cellcolor[rgb]{0.9,0.9,0.9}{
            \makecell[{{p{\linewidth}}}]{
                \texttt{\tiny{[GM$|$GM]}}
                {'word': 'sorcière', 'assignment': 'opponent'}
            }
        }
    }
    & & \\ \\

    \theutterance \stepcounter{utterance}  
    & & & \multicolumn{2}{p{0.3\linewidth}}{
        \cellcolor[rgb]{0.9,0.9,0.9}{
            \makecell[{{p{\linewidth}}}]{
                \texttt{\tiny{[GM$|$GM]}}
                {'hidden': {'team': ['pantalon', 'cour', 'eau', 'stock'], 'innocent': ['arc', 'traverser', 'carré', 'taupe', 'scorpion', 'police', 'triangle', 'avion'], 'opponent': ['bâton', 'pirate', 'tuyau', 'âme'], 'assassin': []}, 'revealed': {'team': {'team': ['papier', 'laser', 'serrure', 'araignée', 'flûte'], 'innocent': [], 'opponent': [], 'assassin': []}, 'opponent': {'team': [], 'innocent': [], 'opponent': ['mercure', 'balle', 'balançoire', 'sorcière'], 'assassin': []}}}
            }
        }
    }
    & & \\ \\

    \theutterance \stepcounter{utterance}  
    & & \multicolumn{4}{p{0.6\linewidth}}{
        \cellcolor[rgb]{0.9,0.9,0.9}{
            \makecell[{{p{\linewidth}}}]{
                \texttt{\tiny{[P1$\langle$GM]}}
                \texttt{Le mot flûte a été deviné correctement. Le tour de votre coéquipier s'est terminé là.} \\
\texttt{Fournissez maintenant un autre indice relatif à certains des mots restants de votre équipe et une liste des mots cibles correspondants. N'oubliez pas de commencer votre indice par 'INDICE: ', de mettre une nouvelle ligne, et de commencer votre liste de mots cibles séparés par des virgules par 'CIBLES: '. Remarque : certains mots ont été supprimés des listes par rapport aux demandes précédentes.} \\
\\ 
\texttt{Les mots restants de votre équipe sont : pantalon, cour, eau, stock.} \\
\texttt{Les mots restants pour votre adversaire sont : bâton, pirate, tuyau, âme, sorcière.} \\
\texttt{Les mots distracteurs restants sont : arc, traverser, carré, taupe, scorpion, police, triangle, avion.} \\
\texttt{Les mots tueurs restants sont : .} \\
            }
        }
    }
    & \\ \\

    \theutterance \stepcounter{utterance}  
    & \multicolumn{4}{p{0.6\linewidth}}{
        \cellcolor[rgb]{0.9,0.9,0.9}{
            \makecell[{{p{\linewidth}}}]{
                \texttt{\tiny{[P1$\rangle$GM]}}
                \texttt{INDICE: vêtement} \\
\texttt{CIBLES: pantalon} \\
            }
        }
    }
    & & \\ \\

    \theutterance \stepcounter{utterance}  
    & & & \multicolumn{2}{p{0.3\linewidth}}{
        \cellcolor[rgb]{0.9,0.9,0.9}{
            \makecell[{{p{\linewidth}}}]{
                \texttt{\tiny{[GM$|$GM]}}
                \texttt{vêtement} \\
            }
        }
    }
    & & \\ \\

    \theutterance \stepcounter{utterance}  
    & & & \multicolumn{2}{p{0.3\linewidth}}{
        \cellcolor[rgb]{0.9,0.9,0.9}{
            \makecell[{{p{\linewidth}}}]{
                \texttt{\tiny{[GM$|$GM]}}
                ['pantalon']
            }
        }
    }
    & & \\ \\

    \theutterance \stepcounter{utterance}  
    & & & \multicolumn{2}{p{0.3\linewidth}}{
        \cellcolor[rgb]{0.9,0.9,0.9}{
            \makecell[{{p{\linewidth}}}]{
                \texttt{\tiny{[GM$|$GM]}}
                {'word': 'pantalon', 'assignment': 'team'}
            }
        }
    }
    & & \\ \\

    \theutterance \stepcounter{utterance}  
    & & \multicolumn{4}{p{0.6\linewidth}}{
        \cellcolor[rgb]{0.9,0.9,0.9}{
            \makecell[{{p{\linewidth}}}]{
                \texttt{\tiny{[GM$\rangle$P2]}}
                \texttt{Le mot flûte a été deviné correctement. Votre tour s'est terminé là.} \\
\texttt{Fournissez maintenant une autre liste, séparée par des virgules, d'au moins 1 et jusqu'à 1 mots de la liste suivante de mots qui se rapportent le mieux ou sont le plus étroitement associés au mot 'vêtement'. N'oubliez pas de commencer votre réponse par DEVINETTE: . Remarque : certains mots ont été supprimés de la liste par rapport aux demandes précédentes.} \\
\\ 
\texttt{pirate, taupe, police, âme, bâton, scorpion, pantalon, avion, stock, tuyau, eau, traverser, triangle, arc, carré, cour} \\
            }
        }
    }
    & \\ \\

    \theutterance \stepcounter{utterance}  
    & & & \multicolumn{4}{p{0.6\linewidth}}{
        \cellcolor[rgb]{0.9,0.9,0.9}{
            \makecell[{{p{\linewidth}}}]{
                \texttt{\tiny{[GM$\langle$P2]}}
                \texttt{DEVINETTE: pantalon} \\
            }
        }
    }
     \\ \\

    \theutterance \stepcounter{utterance}  
    & & & \multicolumn{2}{p{0.3\linewidth}}{
        \cellcolor[rgb]{0.9,0.9,0.9}{
            \makecell[{{p{\linewidth}}}]{
                \texttt{\tiny{[GM$|$GM]}}
                ['pantalon']
            }
        }
    }
    & & \\ \\

    \theutterance \stepcounter{utterance}  
    & & & \multicolumn{2}{p{0.3\linewidth}}{
        \cellcolor[rgb]{0.9,0.9,0.9}{
            \makecell[{{p{\linewidth}}}]{
                \texttt{\tiny{[GM$|$GM]}}
                {'word': 'pantalon', 'assignment': 'team'}
            }
        }
    }
    & & \\ \\

    \theutterance \stepcounter{utterance}  
    & & & \multicolumn{2}{p{0.3\linewidth}}{
        \cellcolor[rgb]{0.9,0.9,0.9}{
            \makecell[{{p{\linewidth}}}]{
                \texttt{\tiny{[GM$|$GM]}}
                {'word': 'pantalon', 'assignment': 'team'}
            }
        }
    }
    & & \\ \\

    \theutterance \stepcounter{utterance}  
    & & & \multicolumn{2}{p{0.3\linewidth}}{
        \cellcolor[rgb]{0.9,0.9,0.9}{
            \makecell[{{p{\linewidth}}}]{
                \texttt{\tiny{[GM$|$GM]}}
                {'word': 'bâton', 'assignment': 'opponent'}
            }
        }
    }
    & & \\ \\

    \theutterance \stepcounter{utterance}  
    & & & \multicolumn{2}{p{0.3\linewidth}}{
        \cellcolor[rgb]{0.9,0.9,0.9}{
            \makecell[{{p{\linewidth}}}]{
                \texttt{\tiny{[GM$|$GM]}}
                {'hidden': {'team': ['cour', 'eau', 'stock'], 'innocent': ['arc', 'traverser', 'carré', 'taupe', 'scorpion', 'police', 'triangle', 'avion'], 'opponent': ['pirate', 'tuyau', 'âme'], 'assassin': []}, 'revealed': {'team': {'team': ['papier', 'laser', 'serrure', 'araignée', 'flûte', 'pantalon'], 'innocent': [], 'opponent': [], 'assassin': []}, 'opponent': {'team': [], 'innocent': [], 'opponent': ['mercure', 'balle', 'balançoire', 'sorcière', 'bâton'], 'assassin': []}}}
            }
        }
    }
    & & \\ \\

    \theutterance \stepcounter{utterance}  
    & & \multicolumn{4}{p{0.6\linewidth}}{
        \cellcolor[rgb]{0.9,0.9,0.9}{
            \makecell[{{p{\linewidth}}}]{
                \texttt{\tiny{[P1$\langle$GM]}}
                \texttt{Le mot pantalon a été deviné correctement. Le tour de votre coéquipier s'est terminé là.} \\
\texttt{Fournissez maintenant un autre indice relatif à certains des mots restants de votre équipe et une liste des mots cibles correspondants. N'oubliez pas de commencer votre indice par 'INDICE: ', de mettre une nouvelle ligne, et de commencer votre liste de mots cibles séparés par des virgules par 'CIBLES: '. Remarque : certains mots ont été supprimés des listes par rapport aux demandes précédentes.} \\
\\ 
\texttt{Les mots restants de votre équipe sont : cour, eau, stock.} \\
\texttt{Les mots restants pour votre adversaire sont : bâton, pirate, tuyau, âme.} \\
\texttt{Les mots distracteurs restants sont : arc, traverser, carré, taupe, scorpion, police, triangle, avion.} \\
\texttt{Les mots tueurs restants sont : .} \\
            }
        }
    }
    & \\ \\

    \theutterance \stepcounter{utterance}  
    & \multicolumn{4}{p{0.6\linewidth}}{
        \cellcolor[rgb]{0.9,0.9,0.9}{
            \makecell[{{p{\linewidth}}}]{
                \texttt{\tiny{[P1$\rangle$GM]}}
                \texttt{INDICE: inventaire} \\
\texttt{CIBLES: stock} \\
            }
        }
    }
    & & \\ \\

    \theutterance \stepcounter{utterance}  
    & & & \multicolumn{2}{p{0.3\linewidth}}{
        \cellcolor[rgb]{0.9,0.9,0.9}{
            \makecell[{{p{\linewidth}}}]{
                \texttt{\tiny{[GM$|$GM]}}
                \texttt{inventaire} \\
            }
        }
    }
    & & \\ \\

    \theutterance \stepcounter{utterance}  
    & & & \multicolumn{2}{p{0.3\linewidth}}{
        \cellcolor[rgb]{0.9,0.9,0.9}{
            \makecell[{{p{\linewidth}}}]{
                \texttt{\tiny{[GM$|$GM]}}
                ['stock']
            }
        }
    }
    & & \\ \\

    \theutterance \stepcounter{utterance}  
    & & & \multicolumn{2}{p{0.3\linewidth}}{
        \cellcolor[rgb]{0.9,0.9,0.9}{
            \makecell[{{p{\linewidth}}}]{
                \texttt{\tiny{[GM$|$GM]}}
                {'word': 'stock', 'assignment': 'team'}
            }
        }
    }
    & & \\ \\

    \theutterance \stepcounter{utterance}  
    & & \multicolumn{4}{p{0.6\linewidth}}{
        \cellcolor[rgb]{0.9,0.9,0.9}{
            \makecell[{{p{\linewidth}}}]{
                \texttt{\tiny{[GM$\rangle$P2]}}
                \texttt{Le mot pantalon a été deviné correctement. Votre tour s'est terminé là.} \\
\texttt{Fournissez maintenant une autre liste, séparée par des virgules, d'au moins 1 et jusqu'à 1 mots de la liste suivante de mots qui se rapportent le mieux ou sont le plus étroitement associés au mot 'inventaire'. N'oubliez pas de commencer votre réponse par DEVINETTE: . Remarque : certains mots ont été supprimés de la liste par rapport aux demandes précédentes.} \\
\\ 
\texttt{pirate, taupe, police, âme, scorpion, avion, stock, tuyau, eau, traverser, triangle, arc, carré, cour} \\
            }
        }
    }
    & \\ \\

    \theutterance \stepcounter{utterance}  
    & & & \multicolumn{4}{p{0.6\linewidth}}{
        \cellcolor[rgb]{0.9,0.9,0.9}{
            \makecell[{{p{\linewidth}}}]{
                \texttt{\tiny{[GM$\langle$P2]}}
                \texttt{DEVINETTE: stock} \\
            }
        }
    }
     \\ \\

    \theutterance \stepcounter{utterance}  
    & & & \multicolumn{2}{p{0.3\linewidth}}{
        \cellcolor[rgb]{0.9,0.9,0.9}{
            \makecell[{{p{\linewidth}}}]{
                \texttt{\tiny{[GM$|$GM]}}
                ['stock']
            }
        }
    }
    & & \\ \\

    \theutterance \stepcounter{utterance}  
    & & & \multicolumn{2}{p{0.3\linewidth}}{
        \cellcolor[rgb]{0.9,0.9,0.9}{
            \makecell[{{p{\linewidth}}}]{
                \texttt{\tiny{[GM$|$GM]}}
                {'word': 'stock', 'assignment': 'team'}
            }
        }
    }
    & & \\ \\

    \theutterance \stepcounter{utterance}  
    & & & \multicolumn{2}{p{0.3\linewidth}}{
        \cellcolor[rgb]{0.9,0.9,0.9}{
            \makecell[{{p{\linewidth}}}]{
                \texttt{\tiny{[GM$|$GM]}}
                {'word': 'stock', 'assignment': 'team'}
            }
        }
    }
    & & \\ \\

    \theutterance \stepcounter{utterance}  
    & & & \multicolumn{2}{p{0.3\linewidth}}{
        \cellcolor[rgb]{0.9,0.9,0.9}{
            \makecell[{{p{\linewidth}}}]{
                \texttt{\tiny{[GM$|$GM]}}
                {'word': 'âme', 'assignment': 'opponent'}
            }
        }
    }
    & & \\ \\

    \theutterance \stepcounter{utterance}  
    & & & \multicolumn{2}{p{0.3\linewidth}}{
        \cellcolor[rgb]{0.9,0.9,0.9}{
            \makecell[{{p{\linewidth}}}]{
                \texttt{\tiny{[GM$|$GM]}}
                {'hidden': {'team': ['cour', 'eau'], 'innocent': ['arc', 'traverser', 'carré', 'taupe', 'scorpion', 'police', 'triangle', 'avion'], 'opponent': ['pirate', 'tuyau'], 'assassin': []}, 'revealed': {'team': {'team': ['papier', 'laser', 'serrure', 'araignée', 'flûte', 'pantalon', 'stock'], 'innocent': [], 'opponent': [], 'assassin': []}, 'opponent': {'team': [], 'innocent': [], 'opponent': ['mercure', 'balle', 'balançoire', 'sorcière', 'bâton', 'âme'], 'assassin': []}}}
            }
        }
    }
    & & \\ \\

    \theutterance \stepcounter{utterance}  
    & & \multicolumn{4}{p{0.6\linewidth}}{
        \cellcolor[rgb]{0.9,0.9,0.9}{
            \makecell[{{p{\linewidth}}}]{
                \texttt{\tiny{[P1$\langle$GM]}}
                \texttt{Le mot stock a été deviné correctement. Le tour de votre coéquipier s'est terminé là.} \\
\texttt{Fournissez maintenant un autre indice relatif à certains des mots restants de votre équipe et une liste des mots cibles correspondants. N'oubliez pas de commencer votre indice par 'INDICE: ', de mettre une nouvelle ligne, et de commencer votre liste de mots cibles séparés par des virgules par 'CIBLES: '. Remarque : certains mots ont été supprimés des listes par rapport aux demandes précédentes.} \\
\\ 
\texttt{Les mots restants de votre équipe sont : cour, eau.} \\
\texttt{Les mots restants pour votre adversaire sont : pirate, tuyau, âme.} \\
\texttt{Les mots distracteurs restants sont : arc, traverser, carré, taupe, scorpion, police, triangle, avion.} \\
\texttt{Les mots tueurs restants sont : .} \\
            }
        }
    }
    & \\ \\

    \theutterance \stepcounter{utterance}  
    & \multicolumn{4}{p{0.6\linewidth}}{
        \cellcolor[rgb]{0.9,0.9,0.9}{
            \makecell[{{p{\linewidth}}}]{
                \texttt{\tiny{[P1$\rangle$GM]}}
                \texttt{INDICE: récréation} \\
\texttt{CIBLES: cour} \\
            }
        }
    }
    & & \\ \\

    \theutterance \stepcounter{utterance}  
    & & & \multicolumn{2}{p{0.3\linewidth}}{
        \cellcolor[rgb]{0.9,0.9,0.9}{
            \makecell[{{p{\linewidth}}}]{
                \texttt{\tiny{[GM$|$GM]}}
                \texttt{récréation} \\
            }
        }
    }
    & & \\ \\

    \theutterance \stepcounter{utterance}  
    & & & \multicolumn{2}{p{0.3\linewidth}}{
        \cellcolor[rgb]{0.9,0.9,0.9}{
            \makecell[{{p{\linewidth}}}]{
                \texttt{\tiny{[GM$|$GM]}}
                ['cour']
            }
        }
    }
    & & \\ \\

    \theutterance \stepcounter{utterance}  
    & & & \multicolumn{2}{p{0.3\linewidth}}{
        \cellcolor[rgb]{0.9,0.9,0.9}{
            \makecell[{{p{\linewidth}}}]{
                \texttt{\tiny{[GM$|$GM]}}
                {'word': 'cour', 'assignment': 'team'}
            }
        }
    }
    & & \\ \\

    \theutterance \stepcounter{utterance}  
    & & \multicolumn{4}{p{0.6\linewidth}}{
        \cellcolor[rgb]{0.9,0.9,0.9}{
            \makecell[{{p{\linewidth}}}]{
                \texttt{\tiny{[GM$\rangle$P2]}}
                \texttt{Le mot stock a été deviné correctement. Votre tour s'est terminé là.} \\
\texttt{Fournissez maintenant une autre liste, séparée par des virgules, d'au moins 1 et jusqu'à 1 mots de la liste suivante de mots qui se rapportent le mieux ou sont le plus étroitement associés au mot 'récréation'. N'oubliez pas de commencer votre réponse par DEVINETTE: . Remarque : certains mots ont été supprimés de la liste par rapport aux demandes précédentes.} \\
\\ 
\texttt{pirate, taupe, police, scorpion, avion, tuyau, eau, traverser, triangle, arc, carré, cour} \\
            }
        }
    }
    & \\ \\

    \theutterance \stepcounter{utterance}  
    & & & \multicolumn{4}{p{0.6\linewidth}}{
        \cellcolor[rgb]{0.9,0.9,0.9}{
            \makecell[{{p{\linewidth}}}]{
                \texttt{\tiny{[GM$\langle$P2]}}
                \texttt{DEVINETTE: cour} \\
            }
        }
    }
     \\ \\

    \theutterance \stepcounter{utterance}  
    & & & \multicolumn{2}{p{0.3\linewidth}}{
        \cellcolor[rgb]{0.9,0.9,0.9}{
            \makecell[{{p{\linewidth}}}]{
                \texttt{\tiny{[GM$|$GM]}}
                ['cour']
            }
        }
    }
    & & \\ \\

    \theutterance \stepcounter{utterance}  
    & & & \multicolumn{2}{p{0.3\linewidth}}{
        \cellcolor[rgb]{0.9,0.9,0.9}{
            \makecell[{{p{\linewidth}}}]{
                \texttt{\tiny{[GM$|$GM]}}
                {'word': 'cour', 'assignment': 'team'}
            }
        }
    }
    & & \\ \\

    \theutterance \stepcounter{utterance}  
    & & & \multicolumn{2}{p{0.3\linewidth}}{
        \cellcolor[rgb]{0.9,0.9,0.9}{
            \makecell[{{p{\linewidth}}}]{
                \texttt{\tiny{[GM$|$GM]}}
                {'word': 'cour', 'assignment': 'team'}
            }
        }
    }
    & & \\ \\

    \theutterance \stepcounter{utterance}  
    & & & \multicolumn{2}{p{0.3\linewidth}}{
        \cellcolor[rgb]{0.9,0.9,0.9}{
            \makecell[{{p{\linewidth}}}]{
                \texttt{\tiny{[GM$|$GM]}}
                {'word': 'tuyau', 'assignment': 'opponent'}
            }
        }
    }
    & & \\ \\

    \theutterance \stepcounter{utterance}  
    & & & \multicolumn{2}{p{0.3\linewidth}}{
        \cellcolor[rgb]{0.9,0.9,0.9}{
            \makecell[{{p{\linewidth}}}]{
                \texttt{\tiny{[GM$|$GM]}}
                {'hidden': {'team': ['eau'], 'innocent': ['arc', 'traverser', 'carré', 'taupe', 'scorpion', 'police', 'triangle', 'avion'], 'opponent': ['pirate'], 'assassin': []}, 'revealed': {'team': {'team': ['papier', 'laser', 'serrure', 'araignée', 'flûte', 'pantalon', 'stock', 'cour'], 'innocent': [], 'opponent': [], 'assassin': []}, 'opponent': {'team': [], 'innocent': [], 'opponent': ['mercure', 'balle', 'balançoire', 'sorcière', 'bâton', 'âme', 'tuyau'], 'assassin': []}}}
            }
        }
    }
    & & \\ \\

    \theutterance \stepcounter{utterance}  
    & & \multicolumn{4}{p{0.6\linewidth}}{
        \cellcolor[rgb]{0.9,0.9,0.9}{
            \makecell[{{p{\linewidth}}}]{
                \texttt{\tiny{[P1$\langle$GM]}}
                \texttt{Le mot cour a été deviné correctement. Le tour de votre coéquipier s'est terminé là.} \\
\texttt{Fournissez maintenant un autre indice relatif à certains des mots restants de votre équipe et une liste des mots cibles correspondants. N'oubliez pas de commencer votre indice par 'INDICE: ', de mettre une nouvelle ligne, et de commencer votre liste de mots cibles séparés par des virgules par 'CIBLES: '. Remarque : certains mots ont été supprimés des listes par rapport aux demandes précédentes.} \\
\\ 
\texttt{Les mots restants de votre équipe sont : eau.} \\
\texttt{Les mots restants pour votre adversaire sont : pirate, tuyau.} \\
\texttt{Les mots distracteurs restants sont : arc, traverser, carré, taupe, scorpion, police, triangle, avion.} \\
\texttt{Les mots tueurs restants sont : .} \\
            }
        }
    }
    & \\ \\

    \theutterance \stepcounter{utterance}  
    & \multicolumn{4}{p{0.6\linewidth}}{
        \cellcolor[rgb]{0.9,0.9,0.9}{
            \makecell[{{p{\linewidth}}}]{
                \texttt{\tiny{[P1$\rangle$GM]}}
                \texttt{INDICE: H2O} \\
\texttt{CIBLES: eau} \\
            }
        }
    }
    & & \\ \\

    \theutterance \stepcounter{utterance}  
    & & & \multicolumn{2}{p{0.3\linewidth}}{
        \cellcolor[rgb]{0.9,0.9,0.9}{
            \makecell[{{p{\linewidth}}}]{
                \texttt{\tiny{[GM$|$GM]}}
                \texttt{h2o} \\
            }
        }
    }
    & & \\ \\

    \theutterance \stepcounter{utterance}  
    & & & \multicolumn{2}{p{0.3\linewidth}}{
        \cellcolor[rgb]{0.9,0.9,0.9}{
            \makecell[{{p{\linewidth}}}]{
                \texttt{\tiny{[GM$|$GM]}}
                ['eau']
            }
        }
    }
    & & \\ \\

    \theutterance \stepcounter{utterance}  
    & & & \multicolumn{2}{p{0.3\linewidth}}{
        \cellcolor[rgb]{0.9,0.9,0.9}{
            \makecell[{{p{\linewidth}}}]{
                \texttt{\tiny{[GM$|$GM]}}
                {'word': 'eau', 'assignment': 'team'}
            }
        }
    }
    & & \\ \\

    \theutterance \stepcounter{utterance}  
    & & \multicolumn{4}{p{0.6\linewidth}}{
        \cellcolor[rgb]{0.9,0.9,0.9}{
            \makecell[{{p{\linewidth}}}]{
                \texttt{\tiny{[GM$\rangle$P2]}}
                \texttt{Le mot cour a été deviné correctement. Votre tour s'est terminé là.} \\
\texttt{Fournissez maintenant une autre liste, séparée par des virgules, d'au moins 1 et jusqu'à 1 mots de la liste suivante de mots qui se rapportent le mieux ou sont le plus étroitement associés au mot 'h2o'. N'oubliez pas de commencer votre réponse par DEVINETTE: . Remarque : certains mots ont été supprimés de la liste par rapport aux demandes précédentes.} \\
\\ 
\texttt{pirate, taupe, police, scorpion, avion, eau, traverser, triangle, arc, carré} \\
            }
        }
    }
    & \\ \\

    \theutterance \stepcounter{utterance}  
    & & & \multicolumn{4}{p{0.6\linewidth}}{
        \cellcolor[rgb]{0.9,0.9,0.9}{
            \makecell[{{p{\linewidth}}}]{
                \texttt{\tiny{[GM$\langle$P2]}}
                \texttt{DEVINETTE: eau} \\
            }
        }
    }
     \\ \\

    \theutterance \stepcounter{utterance}  
    & & & \multicolumn{2}{p{0.3\linewidth}}{
        \cellcolor[rgb]{0.9,0.9,0.9}{
            \makecell[{{p{\linewidth}}}]{
                \texttt{\tiny{[GM$|$GM]}}
                ['eau']
            }
        }
    }
    & & \\ \\

    \theutterance \stepcounter{utterance}  
    & & & \multicolumn{2}{p{0.3\linewidth}}{
        \cellcolor[rgb]{0.9,0.9,0.9}{
            \makecell[{{p{\linewidth}}}]{
                \texttt{\tiny{[GM$|$GM]}}
                {'word': 'eau', 'assignment': 'team'}
            }
        }
    }
    & & \\ \\

    \theutterance \stepcounter{utterance}  
    & & & \multicolumn{2}{p{0.3\linewidth}}{
        \cellcolor[rgb]{0.9,0.9,0.9}{
            \makecell[{{p{\linewidth}}}]{
                \texttt{\tiny{[GM$|$GM]}}
                {'word': 'eau', 'assignment': 'team'}
            }
        }
    }
    & & \\ \\

    \theutterance \stepcounter{utterance}  
    & & & \multicolumn{2}{p{0.3\linewidth}}{
        \cellcolor[rgb]{0.9,0.9,0.9}{
            \makecell[{{p{\linewidth}}}]{
                \texttt{\tiny{[GM$|$GM]}}
                \texttt{team has won} \\
            }
        }
    }
    & & \\ \\

\end{supertabular}
}

\end{document}
